% 来源及取材范围见分题文件；此为整理者转录的正文，不是下载原件。
\begin{lemma}[10.115.7: normalization after localization]
Let $R\to S$ be an injective finite type ring map. Assume $R$ is a domain. Then there exists an integer $d$ and a factorization
\[R\to R[y_1,\ldots,y_d]\to S'\to S\]
by injective maps such that $S'$ is finite over $R[y_1,\ldots,y_d]$ and such that $S'_f\cong S_f$ for some nonzero $f\in R$.
\end{lemma}
\begin{proof}
Pick $x_1,\ldots,x_n\in S$ which generate $S$ over $R$. Let $K$ be the fraction field of $R$ and $S_K=S\otimes_RK$. By Lemma 10.115.4 we can find $y_1,\ldots,y_d\in S$ such that $K[y_1,\ldots,y_d]\to S_K$ is a finite injective map. Note that $y_i\in S$ because we may pick the $y_j$ in the $\mathbf Z$-algebra generated by $x_1,\ldots,x_n$.

As a finite ring map is integral (see Lemma 10.36.3) we can find monic $P_i\in K[y_1,\ldots,y_d][T]$ such that $P_i(x_i)=0$ in $S_K$. Let $f\in R$ be a nonzero element such that $fP_i\in R[y_1,\ldots,y_d][T]$ for all $i$. Then $fP_i(x_i)$ maps to zero in $S_K$. Hence after replacing $f$ by another nonzero element of $R$ we may also assume $fP_i(x_i)$ is zero in $S$.

Set $x'_i=fx_i$ and let $S'\subset S$ be the $R$-subalgebra generated by $y_1,\ldots,y_d$ and $x'_1,\ldots,x'_n$. Note that $x'_i$ is integral over $R[y_1,\ldots,y_d]$ as we have $Q_i(x'_i)=0$ where $Q_i=f^{\deg_T(P_i)}P_i(T/f)$ which is a monic polynomial in $T$ with coefficients in $R[y_1,\ldots,y_d]$ by our choice of $f$. Hence $R[y_1,\ldots,y_d]\subset S'$ is finite by Lemma 10.36.5. Since $S'\subset S$ we have $S'_f\subset S_f$ (localization is exact).

On the other hand, the elements $x_i=x'_i/f$ in $S'_f$ generate $S_f$ over $R_f$ and hence $S'_f\to S_f$ is surjective. Whence $S'_f\cong S_f$ and we win.
\end{proof}

\begin{lemma}[10.118.1: generic freeness over a Noetherian domain]
Let $R\to S$ be a ring map. Let $M$ be an $S$-module. Assume $R$ is Noetherian, $R$ is a domain, $R\to S$ is of finite type, and $M$ is a finite type $S$-module. Then there exists a nonzero $f\in R$ such that $M_f$ is a free $R_f$-module.
\end{lemma}
\begin{proof}
Let $K$ be the fraction field of $R$. Set $S_K=K\otimes_RS$. This is an algebra of finite type over $K$. We will argue by induction on $d=\dim(S_K)$ (which is finite for example by Noether normalization, see Section 10.115). Fix $d\geq0$. Assume we know that the lemma holds in all cases where $\dim(S_K)<d$.

Suppose given $R\to S$ and $M$ as in the lemma with $\dim(S_K)=d$. By Lemma 10.62.1 there exists a filtration $0\subset M_1\subset M_2\subset\ldots\subset M_n=M$ so that $M_i/M_{i-1}$ is isomorphic to $S/\mathfrak q$ for some prime $\mathfrak q$ of $S$. Note that $\dim((S/\mathfrak q)_K)\leq\dim(S_K)$. Also, note that an extension of free modules is free (see basic notion 50). Thus we may assume $M=S$ and that $S$ is a domain of finite type over $R$.

If $R\to S$ has a nontrivial kernel, then take a nonzero $f\in R$ in this kernel. In this case $S_f=0$ and the lemma holds. (This is really the case $d=-1$ and the start of the induction.) Hence we may assume that $R\to S$ is a finite type extension of Noetherian domains.

Apply Lemma 10.115.7 and replace $R$ by $R_f$ (with $f$ as in the lemma) to get a factorization
\[R\subset R[y_1,\ldots,y_d]\subset S\]
where the second extension is finite. Choose $z_1,\ldots,z_r\in S$ which form a basis for the fraction field of $S$ over the fraction field of $R[y_1,\ldots,y_d]$. This gives a short exact sequence
\[0\to R[y_1,\ldots,y_d]^{\oplus r}\xrightarrow{(z_1,\ldots,z_r)}S\to N\to0.\]
By construction $N$ is a finite $R[y_1,\ldots,y_d]$-module whose support does not contain the generic point $(0)$ of $\Spec(R[y_1,\ldots,y_d])$. By Lemma 10.40.5 there exists a nonzero $g\in R[y_1,\ldots,y_d]$ such that $g$ annihilates $N$, so we may view $N$ as a finite module over $S'=R[y_1,\ldots,y_d]/(g)$. Since $\dim(S'_K)<d$ by induction there exists a nonzero $f\in R$ such that $N_f$ is a free $R_f$-module.

Since $(R[y_1,\ldots,y_d])_f\cong R_f[y_1,\ldots,y_d]$ is free also we conclude by the already mentioned fact that an extension of free modules is free.
\end{proof}

\begin{lemma}[10.118.2: generic freeness]
Let $R\to S$ be a ring map. Let $M$ be an $S$-module. Assume $R$ is a domain, $R\to S$ is of finite presentation, and $M$ is an $S$-module of finite presentation. Then there exists a nonzero $f\in R$ such that $M_f$ is a free $R_f$-module.
\end{lemma}
\begin{proof}
Write $S=R[x_1,\ldots,x_n]/(g_1,\ldots,g_m)$. For $g\in R[x_1,\ldots,x_n]$ denote $\overline g$ its image in $S$. We may write $M=S^{\oplus t}/\sum Sn_i$ for some $n_i\in S^{\oplus t}$. Write $n_i=(\overline g_{i1},\ldots,\overline g_{it})$ for some $g_{ij}\in R[x_1,\ldots,x_n]$. Let $R_0\subset R$ be the subring generated by all the coefficients of all the elements $g_i,g_{ij}\in R[x_1,\ldots,x_n]$.

Define $S_0=R_0[x_1,\ldots,x_n]/(g_1,\ldots,g_m)$. Define $M_0=S_0^{\oplus t}/\sum S_0n_i$. Then $R_0$ is a domain of finite type over $\mathbf Z$ and hence Noetherian (see Lemma 10.31.1). Moreover via the injection $R_0\to R$ we have $S\cong R\otimes_{R_0}S_0$ and $M\cong R\otimes_{R_0}M_0$. Applying Lemma 10.118.1 we obtain a nonzero $f\in R_0$ such that $(M_0)_f$ is a free $(R_0)_f$-module. Hence $M_f=R_f\otimes_{(R_0)_f}(M_0)_f$ is a free $R_f$-module.
\end{proof}
